% formal/hk2017-qp-core.tex — HK2017 QP core
% Problem specification and exact per-sigma solver structure.
%
% Part I — Problem, structure, exact algorithm:
%   def:hk2017-problem — full HK2017 capacity problem
%   fact:hk2017 — maximizer = minimum-action simple Reeb orbit
%   rem:min-support — minimum-length maximizing words have beta > 0
%   rem:adjacency-pruning — directed transition prefilter
%   rem:qp-setup — per-sigma notation (A, F, V, beta_0, H', g, gamma_j)
%   prop:basic-properties — A may be empty, F is compact
%   prop:q-on-affine — Q restricted to affine plane
%   prop:critical-points — 3 cases for critical-point structure
%   prop:boundary-vs-interior — B1-B4 (boundary) and I1-I2 (interior)
%   alg:exact-solver — rational-arithmetic per-sigma solver
%   rem:exact-not-compute — what the solver skips (boundary projection)
%   rem:discontinuities — 6 discontinuities for f64
%
% ── Notational conventions ──
%
% Variables and their roles:
%   beta           free/generic variable ("maximize over beta")
%   beta_0         particular solution of C beta = d
%   beta^*         exact optimal beta (maximizer of Q on F)
%   \tilde\beta    f64-computed approximation of beta^*
%   \delta\beta    = \tilde\beta - beta^* (solution error)
%
% Decoration scheme (applies to all quantities, not just beta):
%   plain          exact mathematical value: H, C, V, H', g, gamma_j, Q
%   tilde (~)      f64-computed value: \tilde H, \tilde V, \tilde H', \tilde gamma_j
%   superscript *  optimal / solution: beta^*, Q^* = Q(beta^*)
%   subscript 0    particular / initial: beta_0
%   delta          perturbation / error: delta H = \tilde H - H
%
% Derived quantities (per-sigma, exact):
%   H, C, d        QP matrices (Definition def:hk2017-problem)
%   A              = {beta : C beta = d} (affine constraint set)
%   F              = A cap {beta >= 0} (closed feasible set)
%   V              orthonormal basis for ker(C), so A = beta_0 + im(V)
%   H'             = V^T H V (reduced Hessian on ker(C))
%   g              = V^T H beta_0 (gradient of Q|_A at alpha=0)
%   gamma_j        eigenvalues of H'
%   k              = dim ker(C) = m - rank(C)
%
% Standard reference: Higham, "Accuracy and Stability of Numerical
% Algorithms" (2002).  Plain = exact, tilde = computed follows Higham's
% convention.  Superscript * for optima is standard in optimization.

% input by formal/main.tex — no \documentclass


% ══════════════════════════════════════════════════════════════════════
% Part I: Problem, structure, and exact algorithm
% ══════════════════════════════════════════════════════════════════════
%
% (1) HK2017 full problem
% (2) Simple Reeb orbit structure: minimum-length maximizing words, pruning
% (3) General facts: quadratic forms on affine planes
% (4) Exact per-sigma solver
% (5) Continuity issues for f64


% ── 1. The active-word HK2017 capacity problem ───────────────────────

\begin{definition}[Active-word HK2017 capacity problem]\label{def:hk2017-problem}
Let $K = \{x \in \mathbb{R}^4 : a_i^\top x \leq 1,\; i = 1,\ldots,F\}$
be a convex polytope with dual vertices
$a_1, \ldots, a_F \in \mathbb{R}^4$.
For each subset $S \subseteq \{1, \ldots, F\}$ with $|S| = m \geq 2$
and each cyclic permutation $\sigma$ of~$S$, written in active Reeb traversal
order, define:
\begin{itemize}
  \item \textbf{Action matrix}
    $H(\sigma) \in \mathbb{R}^{m \times m}$:
    for $i < j$,\;
    $H_{ij} = H_{ji} = \omega_0(a_{\sigma(i)},\, a_{\sigma(j)})$;\;
    $H_{ii} = 0$.
    (Symmetric by construction.
    $\omega_0$ is antisymmetric: $\omega_0(u, v) = -\omega_0(v, u)$.
    See Lemma~\ref{lem:H-quadratic}.)
  \item \textbf{Constraint matrix}
    $C(\sigma) \in \mathbb{R}^{5 \times m}$:
    rows $0$--$3$ are $C_{d,k} = a_{\sigma(k),d}$
    (closure: $\sum a_{\sigma(i)} \beta_i = 0$);
    row $4$ is all ones
    (normalization: $\sum \beta_i = 1$).
  \item $d = (0, 0, 0, 0, 1)^\top$.
\end{itemize}
The \emph{active-word HK2017 capacity problem} is:
\begin{equation}\label{eq:hk2017-problem}
  Q_{\max}
  = \max_{\sigma}\;
    \max_{\beta}\;
    \tfrac{1}{2}\,\beta^\top H(\sigma)\, \beta
  \quad \text{s.t.} \quad
  C(\sigma)\beta = d, \;\; \beta \geq 0,
\end{equation}
where $\sigma$ ranges over all cyclic permutations of all subsets
$S \subseteq \{1, \ldots, F\}$ with $|S| \geq 2$.
The EHZ capacity is $c_{\mathrm{EHZ}}(K) = 1/(2\, Q_{\max})$.
This is HK2017 Theorem~1.1 after translating normals/heights to dual vertices
and reversing the ordered word used in HK2017's displayed objective.
\end{definition}

\begin{fact}[Consequence of HK2017]\label{fact:hk2017}
% Citation: HK2017 Theorems 1.1 (formula_theorem: capacity = max of QP),
%   1.5 (simple_loop_theorem: minimizer is simple), and Lemma 2.1 (dual correspondence).
%   Note: HK2017 1.4 is a Remark, not a Theorem.
Any global maximizer $(\sigma^*, \beta^*)$
of~\eqref{eq:hk2017-problem}
corresponds to a minimum-action simple closed Reeb orbit on
$\partial K$, with the word obtained from $\sigma^*$ by deleting its
zero-weight entries in active traversal order.
\end{fact}

\begin{proof}
Delete the zero coordinates of~$\beta^*$ and the corresponding entries of
$\sigma^*$; this does not change feasibility or the objective.  The remaining
coordinates are positive. In the dual-vertex convention of
Definition~\ref{def:hk2017-problem}, put
\[
  Q=Q_{\sigma^*}(\beta^*)=Q_{\max}>0,
  \qquad
  T=\frac{1}{2Q}=c_{\mathrm{EHZ}}(K),
  \qquad
  \tau_k=T\beta_k^*.
\]
Construct a loop~$z$ of period~$T$ whose velocity is
$R_{\sigma^*(k)}=2J_0a_{\sigma^*(k)}$ for time~$\tau_k$. Closure is precisely
$\sum_k\beta_k^*a_{\sigma^*(k)}=0$. The shoelace formula gives
\[
  2A(z)=4T^2Q=2T.
\]
Moreover $h_K(-J_0R_i)=h_K(2a_i)=2$ for every facet direction, so
$I_K(z)=T$. Thus $A(z)=I_K(z)=T=c_{\mathrm{EHZ}}(K)$: $z$ is a global
minimizer of the free-period dual problem, not merely a feasible or stationary
dual curve.

It is therefore weak dual critical. In the free-period Euler--Lagrange
equation, its multiplier is $I_K(z)/T=1$. Clarke's dual action correspondence,
in the nonsmooth form used in \cite[Lemma~2.1]{HK2017}, consequently puts the
loop on~$\partial K$ by translation alone, as a minimum-action generalized
closed characteristic. Translation preserves the ordered pure-velocity
pieces. On a piece with velocity $R_i$, the reconstructed inclusion is
\[
  2a_i=-J_0R_i\in\partial H_K(\gamma)
  =2\operatorname{conv}\{a_j:F_j\text{ is active at }\gamma\}.
\]
Since the normalized presentation is irredundant, $a_i$ is an extreme vertex
of $K^\circ$; hence facet $F_i$ is itself active. Thus, after the zero-duration
pieces have been deleted, the resulting orbit is simple and has the reduced
word obtained from~$\sigma^*$ in active traversal order.
\end{proof}


% ── 2. Structure from Reeb orbits ────────────────────────────────────

\begin{remark}[Minimum-support maximizers]\label{rem:min-support}
A global maximizer $(\sigma^*, \beta^*)$
of~\eqref{eq:hk2017-problem}
may have $\beta^*_k = 0$ for some indices~$k$
(the orbit does not dwell on those facets).
Among all word descriptions of global maximizers achieving~$Q_{\max}$,
take one with minimum word length.
This minimum-length description has $\beta^* > 0$:
if $\beta^*_k = 0$ for some~$k$, removing $\sigma(k)$
from the permutation gives a shorter $\sigma'$ with the same
$Q$ value (the $k$-th row/column of~$H$ contributes nothing,
and the constraints are satisfied by the remaining components).
This contradicts minimality of the word length.

\textbf{Consequence:} the HK2017 algorithm can skip any
$(\sigma, \beta)$ where the maximum of~$Q$ over
$\{C(\sigma)\beta = d,\; \beta \geq 0\}$
is attained only on the boundary ($\beta_k = 0$ for some~$k$).
A shorter~$\sigma'$ achieves the same or higher~$Q$.
\end{remark}

\begin{remark}[Adjacency pruning]\label{rem:adjacency-pruning}
After the minimum-support reduction, one can further skip
permutations~$\sigma$ where consecutive pairs
$(\sigma(i), \sigma(i{+}1))$ fail the directed transition
feasibility test (Lemma~\ref{lem:numerical-transition-feasibility}):
no Reeb trajectory with that combinatorics exists.
This is independent of the $\beta$-optimization and is applied
as a prefilter before solving.
\end{remark}


% ── 3. Quadratic forms on affine planes ──────────────────────────────

% General facts, not specific to EHZ.

\begin{remark}[Setup and notation]\label{rem:qp-setup}
For a fixed~$\sigma$, write $H = H(\sigma)$,
$C = C(\sigma)$ (suppressing the $\sigma$-dependence).
Define:
\begin{itemize}
  \item $\mathcal{A} = \{\beta \in \mathbb{R}^m : C\beta = d\}$
    \quad (affine constraint set).
  \item $\mathcal{F} = \mathcal{A} \cap \{\beta \geq 0\}$
    \quad (closed feasible set).
  \item $\beta_0 \in \mathcal{A}$: a particular solution
    (e.g., $\beta_0 = C^+ d$, the minimum-norm solution).
  \item $V \in \mathbb{R}^{m \times k}$: orthonormal basis for
    $\ker(C)$,\; $k = m - \operatorname{rank}(C)$,
    \;so $\mathcal{A} = \beta_0 + \operatorname{im}(V)$.
  \item $H' = V^\top H V \in \mathbb{R}^{k \times k}$:
    the \emph{reduced Hessian} (restriction of $H$ to~$\mathcal{A}$).
  \item $\gamma_1, \ldots, \gamma_k$: eigenvalues of~$H'$,
    with orthonormal eigenvectors $w_1, \ldots, w_k \in \mathbb{R}^k$.
  \item $g = V^\top H \beta_0 \in \mathbb{R}^k$:
    the gradient of $Q|_{\mathcal{A}}$ at $\alpha = 0$.
\end{itemize}
The per-$\sigma$ problem is:
maximize $Q(\beta) = \tfrac{1}{2}\beta^\top H \beta$
over $\beta \in \mathcal{F}$.
\end{remark}

\begin{proposition}[Basic properties of $\mathcal{A}$ and $\mathcal{F}$]%
\label{prop:basic-properties}
\leavevmode
\begin{enumerate}
  \item $\mathcal{A}$ may be empty (if $C\beta = d$ is inconsistent).
  \item $\mathcal{F} = \mathcal{A} \cap \{\beta \geq 0\}$
    is compact (closed and bounded; contained in the simplex
    $\Delta = \{\beta \geq 0,\; \sum \beta_i = 1\}$
    by the normalization constraint).
    This includes the case $\mathcal{F} = \emptyset$
    (which is compact): even when $\mathcal{A} \neq \emptyset$,
    the intersection with $\{\beta \geq 0\}$ may be empty.
  \item When $\mathcal{F} \neq \emptyset$,
    $Q$ attains its maximum on~$\mathcal{F}$
    (continuous function on a compact set).
\end{enumerate}
\end{proposition}

\begin{proposition}[$Q$ on the affine plane]\label{prop:q-on-affine}
Parametrize $\mathcal{A}$ as $\beta = \beta_0 + V\alpha$,
$\alpha \in \mathbb{R}^k$.  Then
\begin{equation}\label{eq:q-restricted}
  Q(\beta_0 + V\alpha)
  = \underbrace{\tfrac{1}{2}\,\alpha^\top H' \alpha}_{\text{quadratic}}
  + \underbrace{g^\top \alpha}_{\text{linear}}
  + \underbrace{Q(\beta_0)}_{\text{constant}}.
\end{equation}
Setting $\nabla_\alpha = 0$ gives the stationarity condition
$H' \alpha + g = 0$.
\end{proposition}

\begin{proposition}[Critical-point structure]\label{prop:critical-points}
The critical points of $Q|_{\mathcal{A}}$
(solutions of $H'\alpha + g = 0$) form one of three cases:
\begin{enumerate}
  \item \textbf{$H'$ invertible} ($\det H' \neq 0$):
    unique critical point $\alpha^* = -(H')^{-1} g$,
    hence $\beta^* = \beta_0 + V\alpha^*$.
  \item \textbf{$H'$ singular, $g \perp \ker(H')$}
    (i.e., $g \in \operatorname{im}(H')$):
    affine subspace of critical points,
    $\alpha^*_{\mathrm{part}} + \ker(H')$,
    where $\alpha^*_{\mathrm{part}}$ is any particular solution.
    $Q$ takes the same value at all of them
    (Lemma~\ref{lem:well-defined}).
  \item \textbf{$H'$ singular, $g \not\perp \ker(H')$}:
    no critical points.
    $Q|_{\mathcal{A}}$ has no stationary point
    ($H'\alpha + g = 0$ is inconsistent).
\end{enumerate}
\end{proposition}

\begin{proposition}[Boundary vs.\ interior maximum on $\mathcal{F}$]%
\label{prop:boundary-vs-interior}
Write $Q_{\mathcal{F}} = \max_{\beta \in \mathcal{F}} Q(\beta)$
and $\mathcal{F}_{\max} = \arg\max_{\beta \in \mathcal{F}} Q(\beta)
\subseteq \mathcal{F}$.
Define the boundary
$\partial\mathcal{F} = \{\beta \in \mathcal{F} : \beta_k = 0
\text{ for some } k\}$.

\medskip\noindent
\textbf{Cases where}
$\mathcal{F}_{\max} \subseteq \partial\mathcal{F}$
(max is on the boundary):
\begin{enumerate}
  \item[(B1)] Some $\gamma_j > 0$:
    if some $\beta^* \in \mathcal{F}_{\max}$ had $\beta^* > 0$,
    then $\beta^*$ would be a local maximum of $Q|_{\mathcal{A}}$
    (since $\mathcal{F}$ agrees with $\mathcal{A}$ near any
    interior point).
    The second-order necessary condition for a local maximum
    requires $H' \preceq 0$, contradicting $\gamma_j > 0$.
    So $\mathcal{F}_{\max} \subseteq \partial\mathcal{F}$.

  \item[(B2)] All $\gamma_j < 0$ ($H'$ negative definite),
    unique critical point $\beta^* = \beta_0 - V(H')^{-1}g$,
    but $\beta^*_k \leq 0$ for some~$k$:
    $\beta^*$ is the global maximum of~$Q$ on~$\mathcal{A}$
    (strict concavity), and $Q(\beta) < Q(\beta^*)$ for all
    $\beta \in \mathcal{A} \setminus \{\beta^*\}$.
    Since $\beta^* \notin \{\beta > 0\}$,
    $\mathcal{F}_{\max} \subseteq \partial\mathcal{F}$.

  \item[(B3)] $H'$ singular, $g \not\perp \ker(H')$
    (case~3 of Prop.~\ref{prop:critical-points}):
    $Q|_{\mathcal{A}}$ has no critical point
    ($H'\alpha + g = 0$ has no solution).
    If some $\beta^* \in \mathcal{F}_{\max}$ had $\beta^* > 0$,
    then $\beta^*$ would be a local max of $Q|_{\mathcal{A}}$
    (locally $\mathcal{F} = \mathcal{A}$), hence a critical point.
    Contradiction.
    So $\mathcal{F}_{\max} \subseteq \partial\mathcal{F}$.

  \item[(B4)] $H'$ negative semidefinite (all $\gamma_j \leq 0$,
    some $\gamma_j = 0$), $g \perp \ker(H')$,
    critical subspace does not intersect~$\{\beta > 0\}$:
    $Q|_{\mathcal{A}}$ is concave, so its maximum on~$\mathcal{A}$
    is attained on the critical subspace.
    Since no critical point has $\beta > 0$,
    $\mathcal{F}_{\max} \subseteq \partial\mathcal{F}$.
\end{enumerate}

\medskip\noindent
\textbf{Cases where}
$\mathcal{F}_{\max} \cap \{\beta > 0\} \neq \emptyset$
(max is at an interior point):
\begin{enumerate}
  \item[(I1)] All $\gamma_j < 0$ (negative definite $H'$),
    unique critical point $\beta^* > 0$:
    $\beta^*$ is the global maximum on~$\mathcal{A}$
    (concavity) and lies in the interior of~$\mathcal{F}$.
    $Q_{\mathcal{F}} = Q(\beta^*)$.

  \item[(I2)] $H'$ negative semidefinite, $g \perp \ker(H')$,
    critical subspace intersects $\{\beta > 0\}$:
    $Q_{\mathcal{F}} = Q(\beta^*)$ for any critical
    $\beta^* > 0$ in the subspace.
\end{enumerate}

In both~(I1) and~(I2), the interior maximizer has $\beta^* > 0$
and $Q(\beta^*)$ is the per-$\sigma$ contribution to~$Q_{\max}$.
\end{proposition}

\begin{proof}
(B1): If $\mathcal{F}_{\max} \not\subseteq \partial\mathcal{F}$,
some maximizer $\beta^* \in \mathcal{F}$ has $\beta^* > 0$.
Then $\beta^*$ is a local max of $Q$ on~$\mathcal{A}$
(locally $\mathcal{F}$ agrees with $\mathcal{A}$).
The second-order necessary condition requires $H'$ negative
semidefinite, contradicting $\gamma_j > 0$.
(The max may be interior to a \emph{face} of~$\mathcal{F}$
where some $\beta_k = 0$; this is still in~$\partial\mathcal{F}$.)

(B2): $H'$ negative definite means $Q|_{\mathcal{A}}$ is strictly
concave with unique maximum at~$\beta^*$.
For any $\beta \in \mathcal{F}$, $Q(\beta) \leq Q(\beta^*)$.
Since $\beta^* \notin \mathcal{F}^{\circ}$,
the max on~$\mathcal{F}$ is on~$\partial\mathcal{F}$.

(B3): If some $\beta^* \in \mathcal{F}_{\max}$ had $\beta^* > 0$,
it would be a local max of $Q|_{\mathcal{A}}$ (locally $\mathcal{F} = \mathcal{A}$),
hence a critical point. But no critical point exists. Contradiction.

(B4): $Q$ is constant on the critical subspace.
Any $\beta \notin$ critical subspace has strictly lower~$Q$
(the non-null eigenvalues are $< 0$).
Since the critical subspace misses~$\{\beta > 0\}$,
the max on~$\mathcal{F}$ is on~$\partial\mathcal{F}$.

(I1), (I2): The critical $\beta^*$ achieves the global max of~$Q$
on~$\mathcal{A}$ (concavity / negative semidefiniteness).
Since $\beta^* > 0$, it is in~$\mathcal{F}^{\circ}$ and
$Q_{\mathcal{F}} = Q(\beta^*)$.
\end{proof}


% ── 4. Exact per-sigma solver ────────────────────────────────────────

\begin{algorithm}[Exact per-$\sigma$ solver]\label{alg:exact-solver}
Given exact $(H, C, d)$ over~$\mathbb{Q}$:

\medskip\noindent
\textbf{Step~1.}
Compute $\mathcal{A}$, $V$, $\beta_0$, $H'$, $g$, $\gamma_j$
(notation of Remark~\ref{rem:qp-setup}).
If $\mathcal{A} = \emptyset$: return DROP.

\medskip\noindent
\textbf{Step~2 (Boundary test).}
Determine whether $\mathcal{F}_{\max} \subseteq \partial\mathcal{F}$
using Proposition~\ref{prop:boundary-vs-interior}:
\begin{itemize}
  \item Check (B1): any $\gamma_j > 0$?
  \item Check (B3): $H'$ singular and $g \not\perp \ker(H')$?
  \item If $H'$ invertible:
    compute $\beta^* = \beta_0 - V(H')^{-1} g$.
    Check (B2): any $\beta^*_k \leq 0$?
  \item If $H'$ singular and $g \perp \ker(H')$:
    check (B4): does the critical subspace
    miss $\{\beta > 0\}$? (Linear feasibility in~$\gamma$.)
\end{itemize}

If any of (B1)--(B4) holds: return DROP.
By Remark~\ref{rem:min-support}, a shorter
permutation achieves the same~$Q$.
The HK2017 loop finds it.

\medskip\noindent
\textbf{Step~3 (Report).}
We are in case (I1) or (I2).
Return $Q_{\mathcal{F}} = Q(\beta^*)$
for any critical $\beta^*$ with $\beta^* > 0$.
\end{algorithm}

\begin{remark}[What the exact solver does \emph{not} compute]%
\label{rem:exact-not-compute}
When $\mathcal{F}_{\max} \subseteq \partial\mathcal{F}$ (cases B1--B4),
the solver returns DROP without computing $Q_{\mathcal{F}}$ or
$\mathcal{F}_{\max}$.
The complete answer would require projecting onto each boundary
face ($\beta_k = 0$) and optimizing there --- but the HK2017 loop
already does this by enumerating shorter permutations~$\sigma'$
(at least the ones that matter for~$Q_{\max}$).
\end{remark}


% ── 4b. Near-boundary dropping ───────────────────────────────────────

\begin{lemma}[Q loss from dropping a near-boundary component]%
\label{lem:near-boundary-drop}
Let $\beta^* \in \mathcal{F}$ achieve
$Q_{\mathcal{F}} = \max_{\beta \in \mathcal{F}} Q(\beta)$,
with $\beta^*_k \leq \varepsilon$ for some index~$k$
and $\beta^*_i \geq \delta > 0$ for all $i \neq k$.
Let $\sigma' = \sigma$ with $\sigma(k)$ removed,
with constraint matrix $C_{-k} \in \mathbb{R}^{5 \times (m-1)}$
(column~$k$ of~$C$ deleted), and feasible set
$\mathcal{F}' = \{C_{-k}\beta' = d,\; \beta' \geq 0\}$.

Assume $\sigma_{\min}(C_{-k}) > 0$ and
\begin{equation}\label{eq:positivity-condition}
  \delta > \frac{\|c_k\|}{\sigma_{\min}(C_{-k})} \cdot \varepsilon,
\end{equation}
where $c_k = C_{:,k}$ is the $k$-th column of~$C$.
Write $H_{-k} \in \mathbb{R}^{(m-1) \times (m-1)}$ for $H$
with row and column~$k$ deleted,
and $Q_{-k}(\beta') = \tfrac{1}{2}\beta'^\top H_{-k}\, \beta'$.

Then $Q^*(\sigma') \geq Q^*(\sigma) - E_{\mathrm{drop}}$ where
\begin{equation}\label{eq:combined-drop-bound}
  E_{\mathrm{drop}}
  = \varepsilon \left(
    \|H\| \cdot \|\beta^*\|
    + \frac{\|H\| \cdot \|\beta^*\| \cdot \|c_k\|}
           {\sigma_{\min}(C_{-k})}
    + \frac{\|H\| \cdot \|c_k\|^2 \varepsilon}
           {2\, \sigma_{\min}(C_{-k})^2}
  \right).
\end{equation}
\end{lemma}

\begin{proof}
Define $\beta' \in \mathbb{R}^{m-1}$ by deleting component~$k$
from~$\beta^*$ (reindexed to $\{1, \ldots, m-1\}$).

\emph{Step 1: Constraint violation of $\beta'$.}
$C_{-k}\beta' = \sum_{i \neq k} C_{:,i}\, \beta^*_i
= C\beta^* - c_k\, \beta^*_k = d - c_k\, \beta^*_k$.
So $\|C_{-k}\beta' - d\| = \|c_k\| \cdot \beta^*_k
\leq \|c_k\| \cdot \varepsilon$.

\emph{Step 2: Feasibility restoration.}
Define the projection onto $\{C_{-k}\beta' = d\}$:
$\bar\beta' = \beta' + C_{-k}^\top (C_{-k} C_{-k}^\top)^{-1}
  (d - C_{-k}\beta')$.
Then $C_{-k}\bar\beta' = d$ and
\[
  \|\bar\beta' - \beta'\|
  \leq \frac{\|c_k\|}{\sigma_{\min}(C_{-k})} \cdot \beta^*_k
  \leq \frac{\|c_k\|}{\sigma_{\min}(C_{-k})} \cdot \varepsilon.
\]
Componentwise: $\bar\beta'_i \geq \beta^*_i - \|\bar\beta' - \beta'\|
\geq \delta - \|c_k\| \varepsilon / \sigma_{\min}(C_{-k}) > 0$
by~\eqref{eq:positivity-condition}.
So $\bar\beta' \in \mathcal{F}'$ (feasible for~$\sigma'$).

\emph{Step 3: Q change from removing component~$k$.}
Expanding $Q(\beta^*) = \tfrac{1}{2}\sum_{i,j} H_{ij}\beta^*_i\beta^*_j$
and using $H_{kk} = 0$:
\begin{equation}\label{eq:q-removal}
  Q(\beta^*) - Q_{-k}(\beta')
  = \beta^*_k \sum_{j \neq k} H_{kj}\, \beta^*_j.
\end{equation}
Hence $|Q(\beta^*) - Q_{-k}(\beta')|
\leq \beta^*_k \cdot \|H_{k,:}\| \cdot \|\beta^*\|
\leq \varepsilon \cdot \|H\| \cdot \|\beta^*\|$
(Cauchy--Schwarz; $\|H_{k,:}\| \leq \|H\|$).

\emph{Step 4: Q change from the feasibility correction.}
By the triangle inequality on the quadratic Taylor expansion:
\[
  |Q_{-k}(\bar\beta') - Q_{-k}(\beta')|
  \leq \|H_{-k}\| \cdot \|\beta'\| \cdot \|\bar\beta' - \beta'\|
    + \tfrac{1}{2}\|H_{-k}\| \cdot \|\bar\beta' - \beta'\|^2.
\]
Since $\|H_{-k}\| \leq \|H\|$ (principal submatrix)
and $\|\beta'\| \leq \|\beta^*\|$:
\[
  |Q_{-k}(\bar\beta') - Q_{-k}(\beta')|
  \leq \frac{\|H\| \cdot \|\beta^*\| \cdot \|c_k\| \varepsilon}
            {\sigma_{\min}(C_{-k})}
    + \frac{\|H\| \cdot \|c_k\|^2 \varepsilon^2}
            {2\, \sigma_{\min}(C_{-k})^2}.
\]

\emph{Step 5: Combined bound.}
Since $Q^*(\sigma') \geq Q_{-k}(\bar\beta')$ (maximum $\geq$ any feasible value):
\begin{align*}
  Q^*(\sigma) - Q^*(\sigma')
  &\leq Q(\beta^*) - Q_{-k}(\bar\beta') \\
  &\leq |Q(\beta^*) - Q_{-k}(\beta')| + |Q_{-k}(\beta') - Q_{-k}(\bar\beta')|
  \leq E_{\mathrm{drop}}. \qedhere
\end{align*}
\end{proof}


\begin{remark}[Trinary treatment of $\beta > 0$]%
\label{rem:trinary-beta}
The test $\min_k \beta^*_k > 0$ is discontinuous.
The f64 solver computes $\tilde\beta$ and classifies
(with perturbation bound $\eta \geq \|\tilde\beta - \beta^*\|_\infty$):
\begin{itemize}
  \item \textsc{True}
    ($\min_k \tilde\beta_k > \varepsilon_\beta$
    where $\varepsilon_\beta > \eta$):
    $\beta^*_k \geq \tilde\beta_k - \eta > 0$ for all~$k$.
    Certified interior solution.
    $\tilde Q$ is a reliable approximation of $Q_{\mathcal{F}}(\sigma)$
    (the exact per-$\sigma$ optimum on~$\mathcal{F}$).
  \item \textsc{False}
    ($\min_k \tilde\beta_k < -\varepsilon_\beta$):
    $\beta^*_k < 0$ for some~$k$.
    The unconstrained critical point lies outside~$\mathcal{F}$.
    The true $Q_{\mathcal{F}}(\sigma)$ (max over $\beta \geq 0$)
    is lower, and is achieved on~$\partial\mathcal{F}$;
    a shorter~$\sigma'$ finds it.
    Discard this~$\sigma$.
  \item \textsc{Indeterminate}
    (some $\tilde\beta_k \in [-\varepsilon_\beta,\, \varepsilon_\beta]$):
    cannot determine whether $\beta^* > 0$.
    $\tilde Q$ is the unconstrained value $Q(\tilde\beta)$, which
    may \textbf{overestimate} $Q_{\mathcal{F}}(\sigma)$:
    if $\beta^*_k < 0$, the true maximum on~$\mathcal{F}$ is on the
    boundary and lower.
\end{itemize}

\medskip\noindent
\textbf{Indeterminate is not DROP.}
An \textsc{Indeterminate} node cannot be dropped (its
$\tilde Q$ might be an overestimate, so dropping it could
discard a node whose true $Q_{\mathcal{F}}$ is the global max).
Nor can its $\tilde Q$ be trusted as a lower bound on~$Q_{\max}$
(the overestimate means $\tilde Q$ could exceed the true
$Q_{\max}$, producing a capacity that is too small).

\textsc{Indeterminate} must be \textbf{propagated} to the
accumulator, which resolves it lazily:
\begin{enumerate}
  \item Collect all $(\sigma, \tilde Q, \text{verdict})$ triples.
  \item Sort by $\tilde Q$ descending.
  \item The global maximum $Q_{\max}$ is determined by the
    highest $\tilde Q$ among \textsc{True} nodes.
  \item Any \textsc{Indeterminate} node with
    $\tilde Q > Q_{\max}^{\textsc{True}}$
    must be resolved (e.g., by rational arithmetic) to determine
    whether its true $Q_{\mathcal{F}}(\sigma) > Q_{\max}^{\textsc{True}}$.
  \item \textsc{Indeterminate} nodes with
    $\tilde Q \leq Q_{\max}^{\textsc{True}}$ can be ignored
    (even if their true $Q_{\mathcal{F}}$ is slightly different,
    it cannot exceed $Q_{\max}^{\textsc{True}}$ by more than
    $E_{\mathrm{drop}}$ from Lemma~\ref{lem:near-boundary-drop}).
  \item If no \textsc{True} nodes exist, all candidates are
    \textsc{Indeterminate} and must be resolved.
\end{enumerate}

\medskip\noindent
\textbf{Remark on Lemma~\ref{lem:near-boundary-drop}.}
The lemma remains useful for bounding the Q~loss in item~5:
if $\tilde Q \leq Q_{\max}^{\textsc{True}}$ and $\beta^*_k$
is near zero, the true $Q_{\mathcal{F}}(\sigma)$ differs from
the shorter~$\sigma'$ value by at most $E_{\mathrm{drop}}$.
But the lemma does not justify converting \textsc{Indeterminate}
to DROP --- it bounds the loss \emph{if we drop}, not the
overestimate \emph{if we keep}.
\end{remark}


% ── 5. Continuity issues for f64 ────────────────────────────────────

\begin{remark}[Discontinuities in the exact algorithm]%
\label{rem:discontinuities}
Algorithm~\ref{alg:exact-solver} uses several tests that are
\textbf{discontinuous} functions of the input $(a_1, \ldots, a_F)$.
An f64 implementation cannot evaluate these tests exactly.
We list each discontinuity and what makes it problematic.

\begin{enumerate}
  \item \textbf{Sign of $\gamma_j$ (eigenvalues of $H'$).}
    Eigenvalue \emph{values} are continuous functions of~$H'$
    (Weyl's theorem), but $H'$ itself depends on $V$ (the
    null-space basis of~$C$), which changes discontinuously
    when the rank of~$C$ changes.
    Even at fixed rank, the sign test $\gamma_j > 0$ vs.\
    $\gamma_j \leq 0$ is discontinuous at $\gamma_j = 0$.
    Near $\gamma_j = 0$, f64 rounding can flip the sign.

  \item \textbf{Sign of $\beta^*_k$.}
    $\beta^* = \beta_0 - V(H')^{-1} g$ involves $(H')^{-1}$,
    which amplifies errors by $1/|\gamma_j|$ in eigendirection~$j$.
    When $|\gamma_j|$ is small, $\beta^*$ is ill-conditioned:
    small perturbations of $a_i$ cause large changes in~$\beta^*$.
    The test $\beta^*_k > 0$ vs.\ $\beta^*_k \leq 0$ is
    discontinuous at $\beta^*_k = 0$, and near this boundary,
    the ill-conditioning of~$\beta^*$ makes the test unreliable.

  \item \textbf{Rank of $C$.}
    The dimension $k = m - \operatorname{rank}(C)$
    of the null space is an integer-valued function of~$C$.
    A small perturbation can change rank, which changes the
    dimension of~$\mathcal{A}$ and the size of~$H'$.
    In f64, rank is determined by thresholding singular values
    of~$C$, which is discontinuous at the threshold.

  \item \textbf{Rank of $H'$ (semidefinite case).}
    The solvability check $g \perp \ker(H')$
    (Prop.~\ref{prop:critical-points}, case~2 vs.\ case~3)
    depends on the rank of~$H'$.
    Same discontinuity as item~3.

  \item \textbf{Feasibility of the critical subspace (case B4).}
    The LP ``does the critical subspace intersect
    $\{\beta > 0\}$?'' has a binary answer that is discontinuous
    in the LP coefficients, which depend on~$V$, $H'$, $g$.

  \item \textbf{Eigenvalue threshold for the reduced Hessian.}
    An f64 solver must choose a threshold $\varepsilon$ separating
    ``clearly negative'' from ``near-zero'' eigenvalues.
    Eigenvectors with $|\gamma_j| > \varepsilon$ contribute to
    $(H')^{-1}$ (and hence to $\beta^*$); those with
    $|\gamma_j| \leq \varepsilon$ are treated as null
    (and contribute to the search space for $\beta > 0$).
    This inclusion/exclusion is a discontinuous function of
    $\gamma_j$ at $\gamma_j = \pm\varepsilon$.
    At the threshold, the eigenvector's contribution to
    $\beta^*$ is $O(g_j / \varepsilon)$, which can be large.

    % [GAP - Need an argument that at the threshold, the
    % contribution to Q is small (it should be, since Q varies
    % as O(gamma_j * |delta_beta_j|^2) near critical points,
    % and gamma_j = epsilon is small).  But the contribution
    % to the beta > 0 test may NOT be small.]
\end{enumerate}
\end{remark}
